% ─────────────────────────────────────────────────────────────────────────────
% Statistical Significance — NovelGym, matching RAPid-Learn paper style
% ─────────────────────────────────────────────────────────────────────────────

\subsubsection{Statistical Significance}
\label{sec:ng_stat_sig}

To demonstrate that the hybrid agents (RapidLearn family) adapt significantly
faster than the Transfer RL baselines, we perform an unpaired $t$-test~\cite{kim2015t}
with a confidence interval of $95\%$.
The primary metric for comparison is $T_\text{adapt}$, the number of time steps
required to converge after novelty injection, which is the central efficiency
claim of the hybrid architecture.
For each novelty, we compare the best-performing hybrid agent against each
Transfer RL baseline.
Table~\ref{tab:ng_stat_sig} reports the $p$-values and significance verdicts.

\begin{table}[!tb]
\footnotesize
\begin{center}
\begin{tabular}{>{\centering}p{4.4cm}>{\centering}p{1.2cm}>{\centering}p{1.8cm}p{0.1cm}}

\arrayrulecolor{taupegray}\toprule
\textbf{\begin{small}Methods\end{small}} &
{\textbf{\begin{small}$p$-value\end{small}}} &
\textbf{\begin{small}Statistically Significant\end{small}}\tabularnewline
\midrule

% ── AXE ──
\multicolumn{3}{c}{\textbf{Axe}}\tabularnewline
\midrule
RL(PPO) $\leftrightarrow$ TRL+ICM & $0.0097$ & Yes\tabularnewline
RL(PPO) $\leftrightarrow$ TRL     & $<0.001$ & Yes\tabularnewline
\midrule

% ── CHEST ──
\multicolumn{3}{c}{\textbf{Chest}}\tabularnewline
\midrule
\multicolumn{3}{c}{\small Hybrid agents did not detect this novelty — not applicable}\tabularnewline
\midrule

% ── TRADER ──
\multicolumn{3}{c}{\textbf{Trader}}\tabularnewline
\midrule
RL(PPO) $\leftrightarrow$ TRL+ICM & $0.0050$ & Yes\tabularnewline
RL(PPO) $\leftrightarrow$ TRL     & $0.0079$ & Yes\tabularnewline
\midrule

% ── FENCE ──
\multicolumn{3}{c}{\textbf{Fence}}\tabularnewline
\midrule
RL(PPO) $\leftrightarrow$ TRL+ICM & $<0.001$ & Yes\tabularnewline
RL(PPO) $\leftrightarrow$ TRL     & $<0.001$ & Yes\tabularnewline
\midrule

% ── FIRE ──
\multicolumn{3}{c}{\textbf{Fire}}\tabularnewline
\midrule
RL(PPO) $\leftrightarrow$ TRL+ICM & $0.0902$ & No$^\dagger$\tabularnewline
RL(PPO) $\leftrightarrow$ TRL     & $<0.001$ & Yes$^\ddagger$\tabularnewline

\bottomrule
\multicolumn{3}{l}{\small $n=10$ independent seeds per agent; $\alpha = 0.05$.}\\
\multicolumn{3}{l}{\small $^\dagger$ High variance in TRL+ICM ($\pm 198$k steps) limits power.}\\
\multicolumn{3}{l}{\small $^\ddagger$ Significant but reversed: TRL adapts faster than hybrid on Fire.}\\
\end{tabular}
\caption{\small Results of unpaired $t$-test between the best-performing hybrid agent
and each Transfer RL baseline on $T_\text{adapt}$ (time steps to adapt). The hybrid
advantage is statistically significant across Axe, Trader, and Fence novelties.}
\label{tab:ng_stat_sig}
\end{center}
\end{table}

The results confirm that the hybrid architecture adapts significantly faster than
both Transfer RL baselines across the Axe, Trader, and Fence novelties ($p < 0.01$
in all cases).
The Fire novelty presents a notable exception: Transfer RL (without ICM) adapts
significantly \emph{faster} than the hybrid agent ($p < 0.001$), as the
planning component incurs overhead attempting symbolic recovery before delegating
to the RL subproblem.
The comparison against TRL+ICM on Fire is not significant ($p = 0.09$), driven
by the high variance of that agent on this novelty ($\pm 198$k steps).
The Chest novelty is excluded as hybrid agents never triggered adaptation ---
the original plan continued to succeed, leaving the novelty undetected.